\section{Methods}
\subsection{Data, representation and join integrity}
GSE241132 contains three healthy acute-wound donors, labelled PWH26, PWH27 and PWH28, with Skin, Wound1, Wound7 and Wound30 samples \citep{liu2025_wound,geo_gse241132}. Among 65,521 loaded cells, 58,823 joined the author metadata, and 12,259 author-annotated fibroblasts entered the temporal analysis. The fibroblast counts were 2,152, 1,012, 3,821 and 5,274 at the four times, respectively. Annotation is joined using sample identity and original cell-row identity before quality filtering; a filtered positional row number is not a valid substitute. Earlier exploratory outputs based on positional mismatching are invalid and contribute no results here.

The expression encoder is frozen from GSE165816 \citep{theocharidis2022}; the human wound cohort was not used to refit it. It maps panel-normalized counts into a $q=15$ dimensional Gaussian-encoder mean $\mu$. Only 5,383 of 7,002 panel genes are covered in the human wound input (76.9\%); absent panel genes are zero-filled. The flow acts on $\mu\in\R^{15}$, not on simplex topic weights $\theta=\operatorname{softmax}(\mu)$. Successful reconstruction is therefore conditional on this incomplete transferred panel and this coordinate system.

For training-cell set $\mathcal T$, scaling parameters are fitted coordinatewise on $\mathcal T$ only:
\begin{equation}
a_j=\frac1{|\mathcal T|}\sum_{i\in\mathcal T}\mu_{ij},\qquad
b_j=\left[|\mathcal T|^{-1}\sum_{i\in\mathcal T}(\mu_{ij}-a_j)^2\right]^{1/2},\qquad
z_{ij}=\frac{\mu_{ij}-a_j}{\widetilde b_j}.\label{eq:p2scale}
\end{equation}
In Equation~\eqref{eq:p2scale}, $\widetilde b_j=b_j$ for $b_j>0$ and $\widetilde b_j=1$ if $b_j=0$. The identical transformation is applied to source and target evaluation cells. The target is excluded from estimating $a_j,b_j$, and all cells of the held-out donor are excluded in donor transfer. Distances from different training folds are consequently expressed in their respective fold-standardized spaces.

\subsection{Two distinct held-out tasks}
\textbf{Timepoint reconstruction.} All Wound7 cells are withheld from fitting and scaling. Training segments are Skin$\rightarrow$Wound1 and Wound1$\rightarrow$Wound30 within each donor. The fitted field advances the observed Wound1 source to the specified Wound7 coordinate. The pooled analysis uses three donors. Three additional runs omit one donor and evaluate Wound7 in the remaining donors; these are sensitivity analyses, not predictions for the omitted donor.

\textbf{Donor transfer.} A donor is excluded entirely from fitting and scaling; all four time points in the other two donors are retained. Adjacent within-donor segments Skin$\rightarrow$Wound1, Wound1$\rightarrow$Wound7 and Wound7$\rightarrow$Wound30 are used. The held donor's Wound1 cells initialize prediction, while their Wound7 cells are used for evaluation only. Because Wound7 is observed in training donors, this task does not test reconstruction of a globally unobserved time point. The method benchmark and donor-count experiment use this latter design. Table~\ref{tab:p2protocols} lists the exclusions and settings.

\subsection{Time coordinates and conditional flow matching}
Let $P_t^{(d)}$ denote the empirical distribution of standardized fibroblast coordinates for donor $d$ at day $t$. Write $s=g(t)$ for transformed time. Rank time assigns $(0,1/3,2/3,1)$ to days $(0,1,7,30)$, and the other tested transforms are
\begin{equation}
g_{\rm days}(t)=t/30,\qquad
g_{\rm sqrt}(t)=\sqrt{t/30},\qquad
g_{\rm log}(t)=\log(1+t)/\log31.\label{eq:p2times}
\end{equation}
Equation~\eqref{eq:p2times} defines the three day-based alternatives to rank time. For one training donor and segment $(s_a,s_b)$, endpoints are sampled \emph{independently} from the source and target cells of that donor. Thus the training coupling is the empirical product $\pi_{d,ab}=P_{t_a}^{(d)}\otimes P_{t_b}^{(d)}$. It is not an optimal-transport coupling and does not pair the same biological cell across times. With $\lambda\sim U(0,1)$,
\begin{equation}
z_\lambda=(1-\lambda)z_a+\lambda z_b,\quad
s_\lambda=(1-\lambda)s_a+\lambda s_b,\quad
u_{ab}=\frac{z_b-z_a}{s_b-s_a}.\label{eq:p2bridge}
\end{equation}
Equation~\eqref{eq:p2bridge} defines the interpolated state, time and target velocity. The implemented loss cycles over donor--segment pairs and averages squared error over cells and coordinates:
\begin{equation}
\mathcal L_{\rm CFM}(\phi)=\E_{(d,a,b),\pi_{d,ab},\lambda}
\left[\frac1q\left\|v_\phi(z_\lambda,s_\lambda)-u_{ab}\right\|_2^2\right].\label{eq:p2cfm}
\end{equation}
In Equation~\eqref{eq:p2cfm}, the division of the target velocity by $s_b-s_a$ is necessary: the target is velocity per unit transformed time. In particular, compressing the Skin-to-Wound1 segment on the linear-day axis increases its numerical velocity scale. Equations~\eqref{eq:p2bridge}--\eqref{eq:p2cfm} implement conditional flow matching \citep{lipman2023,tong2023,lipman2024guide}. They do not solve the Benamou--Brenier minimum-kinetic-energy problem \citep{benamou2000} or establish globally optimal transport.

Following the neural ordinary-differential-equation formulation \citep{chen2018ode}, predictions solve an initial-value problem and push forward the empirical source measure:
\begin{equation}
\frac{dz}{ds}=v_\phi(z,s),\qquad z(s_a)=z_a,\qquad
\widehat P_{s_b}=(\Phi^\phi_{s_a\rightarrow s_b})_\#\widehat P_{s_a}.\label{eq:p2ode}
\end{equation}
Equation~\eqref{eq:p2ode} defines prediction from the observed source distribution. The network uses two 128-unit SiLU hidden layers and a 32-dimensional sinusoidal time embedding. Training uses Adam with learning rate $10^{-3}$, 8,000 steps and batches of 256. Direct target predictions use 50 fixed fourth-order Runge--Kutta steps. The time-axis diagnostic instead uses 60 steps over the full Wound1-to-Wound30 segment and reads the grid point nearest the transformed Wound7 coordinate; it does not select the grid point with the lowest target distance. A separate fixed-field resolution comparison is described below; finite-step agreement does not prove convergence to an exact solution.

\subsection{What distributional success does and does not identify}
At the ideal population minimizer of the squared loss, the field is a conditional expectation of path velocities, $v^*(z,s)=\E[u_{ab}\mid z_\lambda=z,s_\lambda=s]$, where defined. Under suitable integrability and regularity, the induced marginal path satisfies a continuity equation,
\begin{equation}
\partial_s p_s+\nabla\!\cdot(p_sv^*)=0.\label{eq:p2continuity}
\end{equation}
Equation~\eqref{eq:p2continuity} describes conservation of probability mass in the model. It neither tracks proliferation/death rates nor determines individual-cell ancestry. Multiple couplings and fields can interpolate the same marginals; finite neural optimization adds approximation error. Matching the observed target distribution therefore does not identify a unique biological vector field \citep{lavenant2021}. With only four measurement times, the values of $g$ between observations are analysis choices rather than an identified biological clock.

\subsection{Distance estimand, finite-sample statistic and thresholds}
For independent $X,X'\sim P$ and $Y,Y'\sim Q$ with finite first moments, define energy distance without the optional square root \citep{szekely2013}:
\begin{equation}
D_E(P,Q)=2\E\|X-Y\|_2-\E\|X-X'\|_2-\E\|Y-Y'\|_2.\label{eq:p2ed}
\end{equation}
In Euclidean space, Equation~\eqref{eq:p2ed} is nonnegative and zero exactly when $P=Q$. The implementation uses the V-statistic, including diagonal zero distances:
\begin{align}
\widehat D_E(P_m,Q_n)={}&\frac{2}{mn}\sum_{i=1}^{m}\sum_{j=1}^{n}\|x_i-y_j\|_2
-\frac1{m^2}\sum_{i=1}^{m}\sum_{i'=1}^{m}\|x_i-x_{i'}\|_2\notag\\[-2pt]
&-\frac1{n^2}\sum_{j=1}^{n}\sum_{j'=1}^{n}\|y_j-y_{j'}\|_2.\label{eq:p2edv}
\end{align}
Equation~\eqref{eq:p2edv} is the exact distance between the two empirical measures, not an unbiased U-statistic for the population distance. For baseline $B=\widehat D_E(P_{\rm source},P_{\rm target})>0$ and predicted distance $D$, gain and relative improvement are
\begin{equation}
G=B-D,\qquad I=G/B=1-D/B.\label{eq:p2improvement}
\end{equation}
Equation~\eqref{eq:p2improvement} defines the absolute gain and the fraction reported as a percentage. Stay-still means $P_{\rm prediction}=P_{\rm source}$. Seeded subsampling caps the quadratic distance calculation (1,500 cells per distribution in the initial temporal and donor-transfer runs; the common benchmark uses its recorded 1,200-cell cap). Counts, caps and per-fold values accompany each comparison; independently run analyses need not have identical stochastic point estimates.

The target is split once into seeded disjoint halves, defining $N=\widehat D_E(P_{\rm target,1},P_{\rm target,2})$. The initial timepoint report asks only whether $G>0$, while the time-axis, donor-transfer, method and donor-count comparisons use $G>N$. The latter is a sampling-scale heuristic, not a significance test with known type-I error, a repeated-resampling interval, or an error bound. Selecting the best position on a path using target distances is retrospective; the preassigned Wound7 coordinate supplies the held-out result, and the full scan is a diagnostic rather than a tuned clinical predictor.

\subsection{Donor context and the non-neural predictors}
For a source sample, context concatenates the coordinatewise mean and standard deviation, $c_d=(\bar z_d,\operatorname{sd}(z_d))\in\R^{2q}$. The conditioned field $v_\phi(z,s,c_d)$ uses this context from the training segment's source. At test time, the held donor's source alone supplies $c_d$; its target is not used. Context is zeroed in 20\% of training batches. Both neural fields use the same donor-segment schedule and endpoint-sampling rule, although successive pseudorandom draws need not be identical. Six training-donor evaluations are reported separately from the three held-out-donor folds.

Let $\delta_d=\bar z_{d,7}-\bar z_{d,1}$ for a training donor. Mean displacement predicts $z+\overline\delta$; nearest-donor displacement predicts $z+\delta_{d^*}$, where $d^*$ maximizes cosine similarity between the source contexts. These are translations, not fitted neural fields. The entropic transport baseline \citep{cuturi2013} uses pooled training-source and training-target cells, uniform empirical marginals $a,b$, and normalized squared Euclidean cost $C$:
\begin{equation}
\Gamma_\varepsilon=\arg\min_{\Gamma\geq0,\,\Gamma\mathbf1=a,\,\Gamma^{\mathsf T}\mathbf1=b}
\left\{\langle C,\Gamma\rangle+\varepsilon\sum_{ij}\Gamma_{ij}(\log\Gamma_{ij}-1)\right\},\quad\varepsilon=0.05.\label{eq:p2ot}
\end{equation}
Equation~\eqref{eq:p2ot} defines the entropic transport plan fitted to the training endpoints. Sinkhorn iteration is capped at 2,000 steps, with at most 1,500 cells per endpoint. Row-normalized barycentres define training-source displacements; each held-out source cell receives the average displacement of its 15 nearest training-source neighbours. This is an out-of-sample displacement extension of a transport plan, not transport fitted to the held target.

\subsection{Donor-count estimand and dependence-aware summaries}
For each held donor $d$ and seed $r\in\{0,1,2\}$, let $D_{1,d,r}$ average the errors from the two possible single-training-donor choices, and let $D_{2,d,r}$ be the error with both available training donors. Equal weighting then gives
\begin{equation}
R_k=\frac1{3\cdot3}\sum_{d=1}^{3}\sum_{r=0}^{2}D_{k,d,r},\qquad
\Delta_d=\frac13\sum_{r=0}^{2}(D_{1,d,r}-D_{2,d,r}),\qquad
\overline\Delta=\frac13\sum_d\Delta_d.\label{eq:p2curve}
\end{equation}
Equation~\eqref{eq:p2curve} defines the seed-averaged estimand and donor-block differences. The plotted main curve uses seed 0, with other seeds shown separately. The 4,000-draw percentile interval resamples the three $\Delta_d$ blocks after averaging seeds \citep{efron1993bootstrap}. Overlapping training donors still induce dependence between held-out errors, so even this interval is a descriptive three-block summary, not a population coverage guarantee. Seeds and alternative training subsets are not additional biological replicates. Configuration-specific ranges quantify stochastic instability. Two values of $k$ cannot identify a plateau, a functional learning-curve model or a required number of donors.

\subsection{Geometry and mouse computation}
Donor displacement alignment is $\cos(\delta_d,\delta_e)=\delta_d^{\mathsf T}\delta_e/(\|\delta_d\|\|\delta_e\|)$. The time-versus-donor contrast compares 18 within-donor temporal distances with 12 between-donor same-time distances. Pair-distance bootstrap intervals are descriptive because pairs share donors. The geometric screen used for endpoints $a,b$ and held time $h$ is
\begin{equation}
J=\ind\{\widehat D_E(P_a,P_b)\geq\max[\widehat D_E(P_a,P_h),\widehat D_E(P_h,P_b)]\}.\label{eq:p2geometry}
\end{equation}
Equation~\eqref{eq:p2geometry} is a distance-ordering diagnostic. It does not imply equality in a triangle inequality, identify a geodesic, or establish that a target lies on a unique direct path. We report it as such rather than as proof of interpolability.

GSE326622 is evaluated separately for the source-labelled NDB, PDB and GDB mouse arms \citep{geo_gse326622}. POD7 is excluded from fitting and scaling; POD0, POD2 and POD30 provide the training time points, with POD2 supplying the prediction source. Each mouse fit uses rank time, 6,000 optimization steps and batches of 256. The human panel's mouse-symbol mapping covers 4,182 of 7,002 genes (59.7\%) in each arm, and each arm has one animal per time point. Animal identity and time are therefore confounded. The three arm computations assess protocol behaviour and representation transfer; cell numbers do not supply independent animals or cross-species clinical validation.

\subsection{Numerical-resolution and evaluation-sampling sensitivity}
The saved human expression coordinates and annotations supply a separate seed-0 pooled fit using the stated training exclusions, 8,000 optimization steps and batch size 256. The fitted field is then fixed. Predictions integrate the same 1,012 day-1 cells with 25, 50, 100 and 200 RK4 steps to the prespecified rank-time day-7 coordinate. Distance calculations use one fixed set of 1,500 day-7 cells and double-precision pairwise distances. Root-mean-square coordinate discrepancies are measured relative to the 200-step result (Table~\ref{tab:p2numerics}).

Evaluation sensitivity fixes the 50-step prediction and varies only cell subsampling. At each count of 250, 500 and 1,000 cells per distribution, 50 seeded draws without replacement use common source and target indices for predicted and unchanged-source distances. Each draw additionally samples two disjoint target subsets of that size for the reference distance. We report median improvement, the central 95\% range across draws and the count exceeding the reference (Table~\ref{tab:p2sampling}). These are conditional Monte Carlo summaries. They neither refit the field per draw nor estimate uncertainty across new donors.
